\subsection{Criterios de selección}

Los criterios de selección se definieron en función del objetivo metodológico del trabajo.
Como la investigación requería integrar compresión sin pérdida dentro del backend de
memoria del simulador, se priorizaron aquellos factores que incidían de manera más
directa en la viabilidad de una arquitectura por bloques, en el costo temporal del acceso
a datos comprimidos y en la facilidad de incorporación de la biblioteca en la
implementación experimental.

\begin{table}[H]
    \centering
    \small
    \setlength{\tabcolsep}{4pt}
    \caption{Criterios para la selección de bibliotecas de compresión sin pérdida}
    \label{tab:criterios_bibliotecas}
    \begin{tabular}{>{\raggedright\arraybackslash}p{1.5cm}>{\raggedright\arraybackslash}p{4.8cm}>{\raggedright\arraybackslash}p{0.42\linewidth}>{\centering\arraybackslash}p{1.2cm}}
        \hline
        \textbf{Código} & \textbf{Criterio} & \textbf{Descripción} & \textbf{Peso} \\ \hline
        C1 & Adecuación a almacenamiento comprimido por bloques &
        Evalúa qué tan naturalmente la biblioteca soporta o facilita un esquema de
        datos particionado en bloques independientes, compatible con compresión y
        descompresión local sin necesidad de reconstruir el arreglo completo en cada
        operación. Este criterio recibió la mayor ponderación porque la arquitectura
        propuesta descansa precisamente en ese patrón de acceso. &
        0.35 \\

        C2 & Costo de acceso en la ruta crítica &
        Valora el comportamiento de la biblioteca cuando la compresión y la
        descompresión forman parte de accesos repetidos sobre fragmentos del vector de
        estado. En este contexto, la latencia de recuperación tiene mayor relevancia
        metodológica que una tasa de compresión aislada obtenida fuera de la dinámica
        del simulador. &
        0.30 \\

        C3 & Facilidad de integración en C/C++ sobre buffers en memoria &
        Considera la disponibilidad de una interfaz clara para operar directamente sobre
        buffers residentes en memoria, así como la facilidad de parametrización e
        incorporación en el backend del simulador. Este criterio recoge el esfuerzo de
        integración requerido para construir y evaluar la propuesta experimental. &
        0.20 \\

        C4 & Afinidad con datos numéricos binarios &
        Evalúa la cercanía de la biblioteca con arreglos homogéneos de datos binarios o
        numéricos, en particular de punto flotante, y su compatibilidad con
        transformaciones reversibles que puedan favorecer la compresibilidad sin alterar
        la exactitud numérica. &
        0.15 \\ \hline
    \end{tabular}
\end{table}

La ponderación asignada refleja la prioridad de materializar una arquitectura de
almacenamiento comprimido por bloques dentro de la ruta de ejecución del simulador.
Por esta razón, la adecuación estructural al particionamiento y el costo temporal de
acceso recibieron mayor peso que otros atributos más generales de las bibliotecas de
compresión.
